\chapter{分类与亲缘关系}\label{ch:g6:classification-kinship}

为什么
\cref{ch:g4:classifying-animals} 中的分类方法如此行之有效？把\omterm{def:g1:animals-around-us:animal}{动物}
按\omterm{def:g1:animals-around-us:covering}{羽毛}和\omterm{def:g1:animals-around-us:covering}{皮毛}分类，其他上千种\omterm{def:g6:classification-kinship:attribute}{特征}便
乖乖排好队——\omterm{prop:g3:sorting-living-things:groups}{哺乳动物}都有\omterm{def:g4:classifying-animals:vertebrate}{脊柱}，
有\omterm{def:g1:animals-around-us:covering}{羽毛}的\omterm{def:g1:animals-around-us:animal}{动物}产有壳的卵，无需逐一核对。
\Cref{rem:g4:classifying-animals:tree} 留下的问题
悬而未决。本章以一个词作答——\omterm{prop:g6:classification-kinship:kinship}{亲缘关系}
——并把分类框图重画成一棵家族树。

\section{特征：共有与嵌套}

\begin{definition}[特征]\label{def:g6:classification-kinship:attribute}
称\emph{特征}\index{特征}为\omterm{def:g6:exploring-our-environment:components}{生物}体所
\emph{具有}的、可以核对的表述：
有\omterm{def:g4:classifying-animals:vertebrate}{脊柱}的\omterm{def:g3:bones-and-muscles:skeleton}{骨骼}；有\omterm{def:g1:my-body:parts}{四肢}；有\omterm{def:g1:animals-around-us:covering}{皮毛}且以\omterm{def:g2:animals-grow-up:born}{乳汁}哺育；
有\omterm{def:g1:animals-around-us:covering}{羽毛}。分类建立在特征之上——
\cref{def:g3:sorting-living-things:criterion}中的良好\omterm{def:g3:sorting-living-things:criterion}{标准}，
经过形式化。
\end{definition}

\begin{proposition}[特征相互嵌套]\label{prop:g6:classification-kinship:nest}
在众多\omterm{def:g5:biodiversity:species}{物种}中核对可见，\omterm{def:g6:classification-kinship:attribute}{特征}会归入
\emph{嵌套}的集合——凡有\omterm{def:g1:animals-around-us:covering}{皮毛}的\omterm{def:g1:animals-around-us:animal}{动物}也都有\omterm{def:g1:my-body:parts}{四肢}；
凡有\omterm{def:g1:my-body:parts}{四肢}的\omterm{def:g1:animals-around-us:animal}{动物}也都有\omterm{def:g4:classifying-animals:vertebrate}{脊柱}；绝不会
出现相反的模式（有\omterm{def:g1:animals-around-us:covering}{皮毛}而无\omterm{def:g4:classifying-animals:vertebrate}{脊柱}），也绝不会
出现相交的集合（同时有\omterm{def:g1:animals-around-us:covering}{皮毛}与\omterm{def:g1:animals-around-us:covering}{羽毛}）。这种嵌套
之所以“类群套类群”
（\cref{def:g4:classifying-animals:classification}）
才能契合\omterm{def:g6:exploring-our-environment:components}{生物}界。
\end{proposition}
\admitted

\begin{omfigure}
\begin{tikzpicture}[scale=0.9]
  % nested attribute boxes
  \draw[thick, omDef, fill=omDef!4] (0,0) rectangle (11.4,4.4);
  \node[font=\small\bfseries, omDef] at (2.2,4.1)
    {有脊柱};
  \draw[thick, omMeth, fill=omMeth!5] (0.4,0.35) rectangle (8.6,3.5);
  \node[font=\small\bfseries, omMeth] at (2.35,3.2)
    {有四肢};
  \draw[thick, omProp, fill=omProp!6] (0.8,0.7) rectangle (5.6,2.6);
  \node[font=\small\bfseries, omProp] at (2.6,2.3)
    {有皮毛与乳汁};
  \node[font=\small] at (3.2,1.4) {猫、蝙蝠、鲸};
  \node[font=\small, align=center] at (7.0,1.6)
    {青蛙、蜥蜴，\\ 乌鸫};
  \node[font=\small] at (10.0,1.6) {鳟鱼};
\end{tikzpicture}

{\small \omterm{def:g6:classification-kinship:attribute}{特征}相互嵌套：有\omterm{def:g1:animals-around-us:covering}{皮毛}的\omterm{def:g1:animals-around-us:animal}{动物}位于有\omterm{def:g1:my-body:parts}{四肢}的
\omterm{def:g1:animals-around-us:animal}{动物}之内，有\omterm{def:g1:my-body:parts}{四肢}的\omterm{def:g1:animals-around-us:animal}{动物}位于有\omterm{def:g4:classifying-animals:vertebrate}{脊柱}的\omterm{def:g1:animals-around-us:animal}{动物}之内。鳟鱼
有\omterm{def:g4:classifying-animals:vertebrate}{脊柱}，却位于有\omterm{def:g1:my-body:parts}{四肢}的框之外。}
\end{omfigure}

\begin{example}[三种动物的归位]\label{ex:g6:classification-kinship:placing}
请检查鲸、鳟鱼和乌鸫各自所在的框。
鲸：有\omterm{def:g4:classifying-animals:vertebrate}{脊柱}，有\omterm{def:g1:my-body:parts}{四肢}（前鳍，以及后方隐藏的
残迹），有\omterm{def:g1:animals-around-us:covering}{皮毛}与\omterm{def:g2:animals-grow-up:born}{乳汁}——位于最内层的框，
与猫同处。鳟鱼：有\omterm{def:g4:classifying-animals:vertebrate}{脊柱}，止步——仅在最外层框内。
乌鸫：有\omterm{def:g4:classifying-animals:vertebrate}{脊柱}，有\omterm{def:g1:my-body:parts}{四肢}（两条腿、两只翅膀），
有\omterm{def:g1:animals-around-us:covering}{羽毛}——位于有\omterm{def:g1:my-body:parts}{四肢}的框内，并处于\omterm{def:g1:animals-around-us:covering}{羽毛}自己的
小格中。没有任何\omterm{def:g1:animals-around-us:animal}{动物}会
同时占据两个相互交叉的框。
\end{example}

\section{亲缘关系：嵌套背后的原因}

\begin{proposition}[共有特征意味着共同祖先]\label{prop:g6:classification-kinship:kinship}
嵌套的原因在于遗传
（\cref{rem:g4:animal-reproduction:mix}：幼体从
亲代处获得）。\omterm{def:g5:biodiversity:species}{物种}共有某一\omterm{def:g6:classification-kinship:attribute}{特征}，是因为它们都
\emph{从共同祖先处遗传而来}——
一个遥远的祖先\omterm{def:g5:biodiversity:species}{物种}。两个\omterm{def:g5:biodiversity:species}{物种}共有的\omterm{def:g6:classification-kinship:attribute}{特征}越多，
它们的\omterm{prop:g6:classification-kinship:kinship}{亲缘关系}就越近：
\emph{亲缘关系}\index{亲缘关系}：猫与鲸两者都是
有\omterm{def:g1:animals-around-us:covering}{皮毛}、以\omterm{def:g2:animals-grow-up:born}{乳汁}哺育的\omterm{def:g1:animals-around-us:animal}{动物}，彼此间的\omterm{prop:g6:classification-kinship:kinship}{亲缘关系}
比它们各自与鳟鱼的关系更近。分类类群归根结底是
一个\emph{家族}：祖先及其全部
后代。
\end{proposition}
\admitted

\begin{example}[终于读懂鲸的归类]\label{ex:g6:classification-kinship:whale}
这正是鲸之谜
（\cref{exo:g4:classifying-animals:9}）必须如此解答：
有\omterm{def:g1:animals-around-us:covering}{皮毛}与\omterm{def:g2:animals-grow-up:born}{乳汁}标记着鲸的家系——其祖先是以\omterm{def:g2:animals-grow-up:born}{乳汁}哺育的陆生\omterm{prop:g3:sorting-living-things:groups}{哺乳动物}，
而其鱼样的体形，
是水中生活所致
（\cref{exo:g4:adapted-to-their-world:11}），同样的生活方式把它
烙印在一个亲缘无关的科上。祖先\omterm{def:g6:classification-kinship:attribute}{特征}划分家族；
\omterm{def:g2:where-animals-live:habitat}{栖息地}上的相似之处，可被毫不相干的
陌生者借用。分类学信任前者，
而不信任后者。
\end{example}

\begin{definition}[亲缘树]\label{def:g6:classification-kinship:tree}
\omterm{prop:g6:classification-kinship:kinship}{亲缘关系}用一棵\emph{亲缘树}\index{亲缘树}来表示：
\omterm{def:g1:plants-around-us:tree}{分支}点代表共同祖先，每一\omterm{def:g6:classification-kinship:attribute}{特征}只出现一次，
出现于它初次产生的那一\omterm{def:g1:plants-around-us:tree}{分支}上——
该分叉之后的每个\omterm{def:g5:biodiversity:species}{物种}都遗传它。嵌套框图
与亲缘树承载相同的信息：一个框即是
带着全部小枝的一个\omterm{def:g1:plants-around-us:tree}{分支}。
\end{definition}

\begin{omfigure}
\begin{tikzpicture}[scale=0.95]
  % tree
  \draw[very thick] (0,0) -- (2.2,0);
  % backbone mark
  \fill[omDef] (1.6,0) circle (0.09);
  \node[font=\footnotesize, omDef, align=center] at (1.6,-0.55)
    {脊柱 \\ 出现};
  \draw[very thick] (2.2,0) -- (3.6,1.6) -- (10.4,1.6);
  \draw[very thick] (2.2,0) -- (4.4,-1.4) -- (10.4,-1.4);
  \node[font=\small, right] at (10.5,-1.4) {鳟鱼};
  % four limbs mark
  \fill[omMeth] (3.1,1.05) circle (0.09);
  \node[font=\footnotesize, omMeth, align=center] at (2.4,1.45)
    {四肢 \\ 出现};
  \draw[very thick] (5.0,1.6) -- (6.2,2.8) -- (10.4,2.8);
  % fur mark on upper branch
  \fill[omProp] (5.75,2.35) circle (0.09);
  \node[font=\footnotesize, omProp, align=center] at (5.0,2.9)
    {皮毛与乳汁 \\ 出现};
  \draw[very thick] (6.9,2.8) -- (7.7,3.7) -- (10.4,3.7);
  \node[font=\small, right] at (10.5,3.7) {猫};
  \node[font=\small, right] at (10.5,2.8) {鲸};
  % feathers on middle branch
  \fill[omThm] (7.4,1.6) circle (0.09);
  \node[font=\footnotesize, omThm, align=center] at (7.4,1.05)
    {羽毛 \\ 出现};
  \node[font=\small, right] at (10.5,1.6) {乌鸫};
\end{tikzpicture}

{\small \omterm{prop:g6:classification-kinship:kinship}{亲缘关系}的“口袋树”。每个彩色圆点代表一个
\omterm{def:g6:classification-kinship:attribute}{特征}，只出现一次；其后的每条\omterm{def:g1:plants-around-us:tree}{分支}都继承它。
猫与鲸共享最新的分叉——图中
所示关系最近的亲属。}
\end{omfigure}

\begin{method}[构建一棵小型亲缘关系树]\label{met:g6:classification-kinship:build}
给定少量\omterm{def:g5:biodiversity:species}{物种}：
\begin{enumerate}
  \item 列出各\omterm{def:g5:biodiversity:species}{物种}的\omterm{def:g6:classification-kinship:attribute}{特征}表——\omterm{def:g5:biodiversity:species}{物种}按行排列，可核对的
        \omterm{def:g6:classification-kinship:attribute}{特征}按列排列，具有则打勾；
  \item 按共有的勾号分组，先看覆盖面最广的\omterm{def:g6:classification-kinship:attribute}{特征}：
        嵌套关系即显现于表中；
  \item 画出树形图：每个\omterm{def:g6:classification-kinship:attribute}{特征}对应一个分叉，覆盖最广的
        分叉最靠近树干；每个\omterm{def:g5:biodiversity:species}{物种}恰好位于它所具有
        \omterm{def:g6:classification-kinship:attribute}{特征}对应的各分叉之后；
  \item 读取\omterm{prop:g6:classification-kinship:kinship}{亲缘关系}：两个\omterm{def:g5:biodiversity:species}{物种}的亲近程度由它们
        共有的最后一个分叉决定。
\end{enumerate}
\end{method}

\begin{remark}[树形图悄然蕴含的意义]\label{rem:g6:classification-kinship:implies}
若认真对待这棵树，它将揭示一个极为重大的结论：
沿任意两条\omterm{def:g1:plants-around-us:tree}{分支}往回追溯，它们必定\emph{相遇}。猫与
鳟鱼共享一个祖先\omterm{def:g5:biodiversity:species}{物种}；由此再往前追溯，猫与
蜗牛、雏菊与儿童也必然如此。所有\omterm{def:g6:exploring-our-environment:components}{生物}于是
就都属于同一个家族，其统一性
（见\cref{prop:g5:first-look-at-cells:principle}）是一种家族相似性——而\omterm{def:g5:biodiversity:species}{物种}
且\omterm{def:g5:biodiversity:species}{物种}必须能够\emph{跨代改变}，否则\omterm{def:g1:plants-around-us:tree}{分支}
便永远无法形成。
这一点是否属实，以及它如何运作，
是\cref{ch:g9:evidence-of-evolution}中的
九年级重大问题——这棵树正迫使人们直面它。
\end{remark}

\section{练习}

\begin{exercise}[$\star$]\label{exo:g6:classification-kinship:1}
何为\omterm{def:g6:classification-kinship:attribute}{特征}？给出三个，并写明可核对的表述。
\end{exercise}

\begin{exercise}[$\star$]\label{exo:g6:classification-kinship:2}
\omterm{def:g6:classification-kinship:attribute}{特征}之间\emph{嵌套}是什么意思？给出
本章中三个方框的例子。
\end{exercise}

\begin{exercise}[$\star$]\label{exo:g6:classification-kinship:3}
\omterm{def:g5:biodiversity:species}{物种}为何共有\omterm{def:g6:classification-kinship:attribute}{特征}？依据
\cref{prop:g6:classification-kinship:kinship}作答。
\end{exercise}

\begin{exercise}[$\star$]\label{exo:g6:classification-kinship:4}
在口袋树上，鳟鱼位于哪一个\omterm{def:g6:classification-kinship:attribute}{特征}圆点之后？
鲸位于哪些圆点之后？
\end{exercise}

\begin{exercise}[$\star$]\label{exo:g6:classification-kinship:5}
鲸与猫，还是鲸与鳟鱼，\omterm{prop:g6:classification-kinship:kinship}{亲缘关系}更近？
请由树的分叉读出答案。
\end{exercise}

\begin{exercise}[$\star$]\label{exo:g6:classification-kinship:6}
说明嵌套方框与树枝\omterm{def:g1:plants-around-us:tree}{分支}之间的关系。
\end{exercise}

\begin{exercise}[$\star\star$]\label{exo:g6:classification-kinship:7}
为何分类时相信“\omterm{def:g1:animals-around-us:covering}{皮毛}与\omterm{def:g2:animals-grow-up:born}{乳汁}”，却不信
“外形像鱼”？请利用
\cref{ex:g6:classification-kinship:whale}作答。
\end{exercise}

\begin{exercise}[$\star\star$]\label{exo:g6:classification-kinship:8}
请画出\omterm{def:g6:classification-kinship:attribute}{特征}表，依据
\cref{met:g6:classification-kinship:build}所述方法，列出：鳟鱼、
青蛙、猫、乌鸫——\omterm{def:g6:classification-kinship:attribute}{特征}：\omterm{def:g4:classifying-animals:vertebrate}{脊柱}、\omterm{def:g1:my-body:parts}{四肢}、
\omterm{def:g1:animals-around-us:covering}{皮毛}与\omterm{def:g2:animals-grow-up:born}{乳汁}、\omterm{def:g1:animals-around-us:covering}{羽毛}。
\end{exercise}

\begin{exercise}[$\star\star$]\label{exo:g6:classification-kinship:9}
在口袋树上加入蜥蜴：它将经过哪些
分叉，其\omterm{def:g1:plants-around-us:tree}{分支}从何处分出？
\end{exercise}

\begin{exercise}[$\star\star$]\label{exo:g6:classification-kinship:10}
蝙蝠飞行像乌鸫，哺乳却像猫。
哪一种相似性标志\omterm{prop:g6:classification-kinship:kinship}{亲缘关系}，哪一种
是由生活方式借来的？请依据树的逻辑作出判断。
\end{exercise}

\begin{exercise}[$\star\star$]\label{exo:g6:classification-kinship:11}
“一个类群是一个祖先及其全部后代。”请检验
此句在\omterm{prop:g3:sorting-living-things:groups}{哺乳动物}及
\cref{exo:g4:classifying-animals:11}中的鸭嘴兽上检验：
产卵这一怪异\omterm{def:g6:classification-kinship:attribute}{特征}打破了该类群，还是属于该类群？
\end{exercise}

\begin{exercise}[$\star\star\star$]\label{exo:g6:classification-kinship:12}
沿\cref{rem:g6:classification-kinship:implies}推至其
极端：若每一对\omterm{def:g1:plants-around-us:tree}{分支}确实相遇，则\omterm{def:g5:biodiversity:species}{物种}在漫长时间中必然具有何种性质？
请用一句话写出这一推论 ——
并指出为支持它，你会要求何种
证据。
\end{exercise}

\section{问题：博物馆的标本抽屉}

\begin{problem}[{周末问题——六件标本，一张
特征表，一棵树}]\label{pb:g6:classification-kinship:1}
博物馆借给班级一抽屉标本，共六件：
鳟鱼、青蛙、蜥蜴、鸽子、小鼠，以及
海豚\omterm{def:g3:bones-and-muscles:skeleton}{颅骨}（附鳍的铸型）。任务是按\omterm{def:g6:classification-kinship:attribute}{特征}
将它们分类，并画出\omterm{prop:g6:classification-kinship:kinship}{亲缘关系}树。

\medskip\noindent\textbf{第一部分——表格。}
\begin{enumerate}
  \item 为这六件标本列出\omterm{def:g6:classification-kinship:attribute}{特征}表，\omterm{def:g6:classification-kinship:attribute}{特征}包括：
        \omterm{def:g4:classifying-animals:vertebrate}{脊柱}；\omterm{def:g1:my-body:parts}{四肢}；干燥的鳞片状\omterm{def:g1:the-five-senses:sense}{皮肤}；\omterm{def:g1:animals-around-us:covering}{羽毛}；
        \omterm{def:g1:animals-around-us:covering}{皮毛}和\omterm{def:g2:animals-grow-up:born}{乳汁}。（海豚几乎无毛，
        但博物馆标签证明其幼体以\omterm{def:g2:animals-grow-up:born}{乳汁}哺育——
        \cref{ex:g3:sorting-living-things:surprises}
        解释了为何这一证据足以定论。）
  \item 六件标本都具有哪一项\omterm{def:g6:classification-kinship:attribute}{特征}？
        这种最广泛的共有在\omterm{prop:g6:classification-kinship:kinship}{亲缘关系}上意味着什么？
  \item 哪些标本在“\omterm{def:g1:my-body:parts}{四肢}”一项上打了勾？请解释
        鸽子和海豚的勾选——按照树的逻辑，
        翅膀和鳍属于什么？
  \item 鳟鱼和海豚都生活在水中，为何表格里
        没有“生活在水中”这一\omterm{def:g6:classification-kinship:attribute}{特征}？
        请引用本章的规则说明。
\end{enumerate}

\medskip\noindent\textbf{第二部分——树。}
\begin{enumerate}[resume]
  \item 画出树：分叉的顺序依次为\omterm{def:g4:classifying-animals:vertebrate}{脊柱}、\omterm{def:g1:my-body:parts}{四肢}，
        然后是体被（鳞片、\omterm{def:g1:animals-around-us:covering}{羽毛}、\omterm{def:g1:animals-around-us:covering}{皮毛}），各自位于独立的\omterm{def:g1:plants-around-us:tree}{分支}上，
        并将六件标本全部放入。
  \item 根据你的树，读出抽屉中与小鼠\omterm{prop:g6:classification-kinship:kinship}{亲缘关系}最近的标本，
        以及与鳟鱼\omterm{prop:g6:classification-kinship:kinship}{亲缘关系}最近的标本。
  \item 青蛙已越过“\omterm{def:g1:my-body:parts}{四肢}”一分叉，却排在
        所有体被分叉之前。它裸露湿润的\omterm{def:g1:the-five-senses:sense}{皮肤}
        （见\cref{rem:g4:classifying-animals:confusion}）表明它属于什么，
        它的\omterm{def:g1:plants-around-us:tree}{分支}在哪里分出？
  \item 第七件标本的标签脱落了：它具有干燥的
        鳞片状\omterm{def:g1:the-five-senses:sense}{皮肤}、\omterm{def:g1:my-body:parts}{四肢}和\omterm{def:g4:classifying-animals:vertebrate}{脊柱}。请将它放到树上，
        并在不看标签的情况下说出它的纲。
\end{enumerate}

\medskip\noindent\textbf{第三部分——意义。}
\begin{enumerate}[resume]
  \item 海豚与鳟鱼都具有流线型体形；海豚与小鼠
        则共享\omterm{def:g1:animals-around-us:covering}{皮毛}与\omterm{def:g2:animals-grow-up:born}{乳汁}的祖先\omterm{def:g6:classification-kinship:attribute}{特征}。哪一对在树上
        距离更近？
        为什么树必须忽略流线型这一\omterm{def:g6:classification-kinship:attribute}{特征}？
  \item 说明“\omterm{def:g1:my-body:parts}{四肢}”这一分叉在历史上主张了什么：
        很久以前，哪一个单一的事物，
        青蛙、蜥蜴、鸽子、小鼠和海豚
        都源自它？
  \item 班级的树使每一条\omterm{def:g1:plants-around-us:tree}{分支}
        最终都汇聚到树干上。用一句话说明
        这幅图悄然断言了关于六件标本的什么，
        以及关于
        绘制它的班级成员的什么。
  \item 一位怀疑者说：“你的树不过是一个归档系统。”
        请给出两部分的回应：
        嵌套的\omterm{def:g6:classification-kinship:attribute}{特征}所反映的是一个\emph{事实}，而
        树\emph{解释}了单纯归档
        本会无法解释的什么？
\end{enumerate}
\end{problem}
